\chapter{Introduction to Time-Series Databases}\label{chap:benjamin-introduction}

The amounts of data increase exponentially. Among the most used databases are
Oracle and Mysql. Those databases are great at storing large amounts of data.
What they lack is the ability to retrieve large amounts of data spanning long
time periods. Time-series databases fill this gap. They are optimized to store
and query large amounts of data over a long period of time.

\section{Why Time-Series Databases?}\label{why-time-series-databases}

Conventional databases like Postgres are commonly used to handle large amounts
of user data, such as chats, accounts, and relationships. For even larger
datasets, No-SQL databases such as MongoDB are often horizontally scaled to
accommodate nearly unlimited growth.

However, consider a scenario where logs are collected, stored and analyzed.

\begin{itemize}
  \item
        1000 sensors
  \item
        every 1 second a new observation
  \item
        one year has \(31536000\) seconds
\end{itemize}

% \[
% \text{1000} \cdot 31536000 = 31536000000
% \]

That amounts to 31 billion observations for just one year. Even if this data
does not need that much space, it can still lead to performance issues when
querying it. Time-series databases can have a serious performance boost in this
scenario, especially if the requested data is filtered over time.

Time series databases are databases that are specialized and optimized for time
stamps and data in a big time frame. For example, TimescaleDB, a Postgres
extension, is able to improve performance to up to 350 times faster queries,
44\% faster ingests, and 95\% storage savings with time-series data.~\cite{TimescaleDB:01}

\section{Use Cases of Time-Series Databases}\label{use-cases-of-time-series-databases}

Time-series data are a sequence of successive points in time. This could be any
record with a time stamp. Common examples are as follows.

\begin{itemize}
  \item
        IoT sensors
  \item
        stock marked information
  \item
        logs, error rates or request counts
  \item
        weather patterns/climate data
  \item
        sales or demand forecasting
\end{itemize}

Time-series databases provide advantages that become most apparent when working
with large datasets.